% =====================================================================
% Round 3: Mathematical Verification Report
% =====================================================================
%
% This document reports the results of an independent, computation-by-
% computation verification of every numerical example in round2_examples.tex
% and every proof in round2_math.tex.  Errors are flagged with severity
% ratings: CRITICAL (wrong final answer), ERROR (wrong intermediate value
% that could mislead a student), or NOTE (cosmetic or rounding issue).
%
% Structure:
%   Part A — Numerical verification of the six examples
%   Part B — Logical audit of the four proofs
%   Part C — Consistency with lecture_notes.tex (macros, labels, refs)
%   Part D — Corrected LaTeX for all items requiring changes
%
% =====================================================================


% =====================================================================
%                     PART A: NUMERICAL VERIFICATION
% =====================================================================

\section*{Part A: Numerical Verification of Examples}

% ---------------------------------------------------------------------
% EXAMPLE 1: "Your First Regression in R^3"
% Verdict: ALL CORRECT
% ---------------------------------------------------------------------

\subsection*{Example 1 (Your First Regression in $\R^3$)}

\textbf{Verdict: ALL CORRECT.}  Every matrix product, inverse, and
arithmetic operation was recomputed from scratch.

\begin{itemize}[nosep]
  \item $\bX^\top\bX = \bigl(\begin{smallmatrix}3&6\\6&14\end{smallmatrix}\bigr)$. \checkmark
  \item $\bX^\top\by = (6,\;13)^\top$. \checkmark
  \item $\det(\bX^\top\bX) = 42-36 = 6$. \checkmark
  \item $(\bX^\top\bX)^{-1} = \bigl(\begin{smallmatrix}7/3&-1\\-1&1/2\end{smallmatrix}\bigr)$. \checkmark
  \item $\bhat = (1,\;1/2)^\top$: intermediate products $14-13=1$ and $-6+13/2 = 1/2$. \checkmark
  \item $\yhat = (3/2,\;2,\;5/2)^\top$. \checkmark
  \item $\be = (-1/2,\;1,\;-1/2)^\top$. \checkmark
  \item $\bX^\top\be = (0,\;0)^\top$: first component $-1/2+1-1/2=0$, second $-1/2+2-3/2=0$. \checkmark
  \item $\bX(\bX^\top\bX)^{-1} = \bigl(\begin{smallmatrix}4/3&-1/2\\1/3&0\\-2/3&1/2\end{smallmatrix}\bigr)$. \checkmark
  \item $\bH = \frac{1}{6}\bigl(\begin{smallmatrix}5&2&-1\\2&2&2\\-1&2&5\end{smallmatrix}\bigr)$. \checkmark
  \item $\tr(\bH) = 12/6 = 2$. \checkmark
  \item $\bH^2$: the $6\times 6$ entry-by-entry product
        $\frac{1}{36}\bigl(\begin{smallmatrix}30&12&-6\\12&12&12\\-6&12&30\end{smallmatrix}\bigr)
        = \frac{1}{6}\bigl(\begin{smallmatrix}5&2&-1\\2&2&2\\-1&2&5\end{smallmatrix}\bigr) = \bH$. \checkmark
\end{itemize}


% ---------------------------------------------------------------------
% EXAMPLE 2: "Pythagoras Checks Out"
% Verdict: ALL CORRECT
% ---------------------------------------------------------------------

\subsection*{Example 2 (Pythagoras Checks Out)}

\textbf{Verdict: ALL CORRECT.}

\begin{itemize}[nosep]
  \item $\bar{y} = 6/3 = 2$. \checkmark
  \item SST $= 1+1+0 = 2$. \checkmark
  \item SSR $= 1/4+0+1/4 = 1/2$. \checkmark
  \item SSE $= 1/4+1+1/4 = 3/2$. \checkmark
  \item SSR $+$ SSE $= 1/2+3/2 = 2 =$ SST. \checkmark
  \item $R^2 = (1/2)/2 = 1/4 = 0.25$. \checkmark
  \item $\cos\theta = 1/2$, $\theta = 60°$. \checkmark
\end{itemize}


% ---------------------------------------------------------------------
% EXAMPLE 3: "FWL by Hand"
% Verdict: FINAL ANSWER CORRECT, but intermediate Step B3 has an ERROR
% ---------------------------------------------------------------------

\subsection*{Example 3 (FWL by Hand)}

\textbf{Verdict: Final answer $\hat\beta_2^{\text{FWL}} = 3$ is correct,
but Step~B3 contains an arithmetic error that propagates through two
subsequent displayed equations.}

\medskip
\noindent\textbf{Part A (full regression): correct.}
\begin{itemize}[nosep]
  \item $\bX^\top\bX = \bigl(\begin{smallmatrix}4&2&2\\2&2&1\\2&1&2\end{smallmatrix}\bigr)$. \checkmark
  \item $\bX^\top\by = (12,\;8,\;9)^\top$. \checkmark
  \item Row reduction verified: $\hat\beta_2=3$, $\hat\beta_1=2$, $\hat\beta_0=1/2$. \checkmark
  \item $\yhat = (5/2,\;7/2,\;11/2,\;1/2)^\top$, $\be = (-1/2,\;-1/2,\;1/2,\;1/2)^\top$. \checkmark
  \item $\bX^\top\be = (0,0,0)^\top$. \checkmark
\end{itemize}

\medskip
\noindent\textbf{Part B (FWL two-step): Steps B1--B2 correct, Step B3 has error.}
\begin{itemize}[nosep]
  \item Step B1: $(\bX_1^\top\bX_1)^{-1} = \bigl(\begin{smallmatrix}1/2&-1/2\\-1/2&1\end{smallmatrix}\bigr)$. \checkmark
  \item Step B2: $\tilde{\bm{x}}_2 = (-1/2,\;1/2,\;1/2,\;-1/2)^\top$. \checkmark
  \item[\bfseries !] \textbf{Step B3 --- ERROR:}
        The document computes
        \[
          (\bX_1^\top\bX_1)^{-1}\bX_1^\top\by
          = \begin{pmatrix}1/2&-1/2\\-1/2&1\end{pmatrix}
            \begin{pmatrix}12\\8\end{pmatrix}
          = \begin{pmatrix}2\\4\end{pmatrix}.
        \]
        The correct computation is:
        \[
          \begin{pmatrix}(1/2)(12)+(-1/2)(8)\\(-1/2)(12)+(1)(8)\end{pmatrix}
          = \begin{pmatrix}6-4\\-6+8\end{pmatrix}
          = \begin{pmatrix}2\\2\end{pmatrix}.
        \]
        The document writes $(2,\;4)^\top$ but the correct answer is
        $\boxed{(2,\;2)^\top}$.

  \item[\bfseries !] \textbf{Cascading errors in Step B3:}
        \begin{itemize}[nosep]
          \item The projection $\bX_1(2,\;4)^\top$ is given as $(6,\;2,\;6,\;2)^\top$.
                With the correct coefficient $(2,\;2)^\top$, the projection is
                $\bX_1(2,\;2)^\top = (2+2,\;2+0,\;2+2,\;2+0)^\top
                = \boxed{(4,\;2,\;4,\;2)^\top}$.
          \item The residualised response is given as
                $\tilde{\by} = (-4,\;1,\;0,\;-1)^\top$.
                The correct value is
                $\tilde{\by} = (2-4,\;3-2,\;6-4,\;1-2)^\top
                = \boxed{(-2,\;1,\;2,\;-1)^\top}$.
        \end{itemize}

  \item Step B4: The final regression $\hat\beta_2 = \tilde{\bm{x}}_2^\top\tilde{\by}\,/\,
        \tilde{\bm{x}}_2^\top\tilde{\bm{x}}_2$ gives $3/1 = 3$ with \emph{either}
        version of $\tilde{\by}$, because the error happens to cancel in the
        inner product.  (With the wrong $\tilde{\by}$:
        $(-1/2)(-4)+(1/2)(1)+(1/2)(0)+(-1/2)(-1) = 2+1/2+0+1/2 = 3$.
        With the correct $\tilde{\by}$:
        $(-1/2)(-2)+(1/2)(1)+(1/2)(2)+(-1/2)(-1) = 1+1/2+1+1/2 = 3$.)
        So the \emph{final answer} $\hat\beta_2 = 3$ is correct despite the
        intermediate error.
\end{itemize}

\medskip
\noindent\textbf{Severity: ERROR.}
A student following along would obtain $\tilde{\by} = (-2,1,2,-1)^\top$
and be confused by the disagreement with the printed $(- 4,1,0,-1)^\top$.
The corrected LaTeX is provided in Part~D below.


% ---------------------------------------------------------------------
% EXAMPLE 4: "Multicollinearity in Action"
% Verdict: Case A correct; Case B has an approximate-number error
% ---------------------------------------------------------------------

\subsection*{Example 4 (Multicollinearity in Action)}

\textbf{Verdict: Case A entirely correct.  Case B has an error in the
approximate coefficient values after perturbation.}

\medskip
\noindent\textbf{Case A: correct.}
\begin{itemize}[nosep]
  \item Sample correlation $r_{12} = 0$: verified from the dot product of centered predictors. \checkmark
  \item Perturbed $\bX^\top\by_{\text{new}} = (13,\;8,\;9)^\top$. \checkmark
  \item Row reduction gives $\hat\beta_1^{\text{new}} = 3/2$, $\hat\beta_2^{\text{new}} = 5/2$,
        $\hat\beta_0^{\text{new}} = 5/4$. \checkmark
  \item $|\Delta\hat\beta_1| = 1/2$, $|\Delta\hat\beta_2| = 1/2$. \checkmark
\end{itemize}

\medskip
\noindent\textbf{Case B: error in approximate values.}
\begin{itemize}[nosep]
  \item $\bX^\top\bX$ with $x_{2,4}=3.01$: verified entry by entry. \checkmark
  \item Centered cross-product matrix $C$: $C_{11}=5$, $C_{12}=5.015$, $C_{22}=5.030075$. \checkmark
  \item $\det(C) = 0.00015$: verified. \checkmark
        (The document's initial ``$-0.0001$'' is explicitly self-corrected in the text; this is acceptable.)
  \item Eigenvalues $\approx 10.03$ and $\approx 0.000015$: verified. \checkmark
  \item Condition number $\approx 670{,}000$: verified ($670{,}681$). \checkmark
  \item[\bfseries !] \textbf{ERROR in perturbed solution:}
        The document claims that changing $y_4$ from $9$ to $10$ yields
        $\bhat \approx (3,\;-65,\;68)$ with a swing of $\pm 67$.
        The actual solution (verified by direct computation of
        $(\bX^\top\bX)^{-1}\bX^\top\by$) is
        $\bhat \approx \boxed{(3,\;-98,\;100)}$
        with a swing of $\pm 100$.
  \item[\bfseries !] \textbf{Cascading error in the summary table:}
        The table entry ``$\Delta\hat\beta$ per unit $\Delta y$: $\approx 67$''
        should be $\boxed{\approx 100}$.
\end{itemize}

\medskip
\noindent\textbf{Severity: ERROR.}
The qualitative message (``the swing is huge'') is correct, but the
specific numbers are wrong by roughly a factor of 1.5.  A student who
solves the system will get $\approx\pm 100$ and be confused.
Corrected LaTeX is provided in Part~D.


% ---------------------------------------------------------------------
% EXAMPLE 5: "Leverage and Influence"
% Verdict: ALL CORRECT (values use legitimate rounding)
% ---------------------------------------------------------------------

\subsection*{Example 5 (Leverage and Influence)}

\textbf{Verdict: ALL CORRECT.}

\begin{itemize}[nosep]
  \item $\bX^\top\bX = \bigl(\begin{smallmatrix}5&20\\20&130\end{smallmatrix}\bigr)$, $\det = 250$. \checkmark
  \item $(\bX^\top\bX)^{-1} = \frac{1}{250}\bigl(\begin{smallmatrix}130&-20\\-20&5\end{smallmatrix}\bigr)$. \checkmark
  \item Leverage formula $h_{ii} = (26-8x_i+x_i^2)/50$: verified for all five observations. \checkmark
  \item $\sum h_{ii} = 100/50 = 2 = p$. \checkmark
  \item $\bhat = (36/25,\;47/50)^\top = (1.44,\;0.94)^\top$: exact fractions. \checkmark
  \item Fitted values and residuals: all correct to the displayed precision (which is exact). \checkmark
  \item SSE $= 131/50 = 2.62$, $\hat\sigma^2 = 131/150 \approx 0.873$. \checkmark
  \item Influence experiment ($y_5$ shift): $\Delta\bhat = (-1.4,\;0.6)^\top$. \checkmark
  \item Influence experiment ($y_3$ shift): $\Delta\bhat = (1.4,\;-0.1)^\top$. \checkmark
  \item Ratio $|\Delta\hat\beta_1| = 0.6/0.1 = 6$. \checkmark
  \item Cook's distances: verified for all five observations, matching the
        tabulated values to two decimal places. \checkmark
\end{itemize}

\medskip
\noindent\textbf{Minor note.}
The table states $p\hat\sigma^2 = 2\times 0.873 = 1.746$.
The exact value is $2 \times 131/150 = 262/150 \approx 1.7467$.
The document rounds $\hat\sigma^2$ to $0.873$ and then multiplies;
$2\times 0.873 = 1.746$ is consistent with that rounding.
No correction needed.


% ---------------------------------------------------------------------
% EXAMPLE 6: "Ridge Shrinkage Direction by Direction"
% Verdict: ALL CORRECT
% ---------------------------------------------------------------------

\subsection*{Example 6 (Ridge Shrinkage)}

\textbf{Verdict: ALL CORRECT.}

\begin{itemize}[nosep]
  \item Eigendecomposition $\bX^\top\bX = \bQ\bD\bQ^\top = \bigl(\begin{smallmatrix}5&4\\4&5\end{smallmatrix}\bigr)$:
        verified. \checkmark
  \item $\bX^\top\by = (19,\;17)^\top$ from $\bX^\top\bX\,\bhat_{\text{OLS}}$. \checkmark
  \item Eigenbasis decomposition: $\alpha_1 = 2$, $\alpha_2 = 1$. \checkmark
  \item Shrinkage formula $\alpha_k^{\text{R}} = \frac{\lambda_k}{\lambda_k+\lambda}\alpha_k^{\text{OLS}}$:
        verified for all four $\lambda$ values. \checkmark
  \item Coefficient reconstruction: $\hat\beta_1 = \alpha_1+\alpha_2$,
        $\hat\beta_2 = \alpha_1-\alpha_2$: verified. \checkmark
  \item Table entries for $\lambda=0,1$: exact. \checkmark
  \item Table entries for $\lambda=10$: $\hat\beta_1 = 217/209 \approx 1.04$,
        $\hat\beta_2 = 179/209 \approx 0.86$. Rounding is legitimate. \checkmark
  \item Table entries for $\lambda=100$: $\hat\beta_1 \approx 0.175$,
        $\hat\beta_2 \approx 0.155$. The document rounds to $0.18$ and $0.16$.
        Slightly generous rounding but acceptable for a table. \checkmark
  \item Direct verification for $\lambda=1$:
        $(\bX^\top\bX+\bI)^{-1} = \frac{1}{20}\bigl(\begin{smallmatrix}6&-4\\-4&6\end{smallmatrix}\bigr)$,
        $\bhat_{\text{Ridge}} = (23/10,\;13/10)^\top$. \checkmark
  \item Norm column: $\|\bhat\| \approx 3.16, 2.64, 1.35$ for $\lambda = 0,1,10$. \checkmark
  \item Norm for $\lambda=100$: exact value is $\approx 0.234$; document rounds to $0.24$.
        Acceptable. \checkmark
\end{itemize}


% =====================================================================
%                     PART B: PROOF VERIFICATION
% =====================================================================

\section*{Part B: Logical Audit of Proofs}


% ----- Projection Theorem -----
\subsection*{Proof of the Projection Theorem}

\textbf{Verdict: CORRECT and COMPLETE.}

The proof has three steps (existence, uniqueness, orthogonality) plus a
converse.  All are logically sound:

\begin{itemize}[nosep]
  \item \textbf{Step 1 (Existence):}  Uses compactness of a closed bounded
        subset of a finite-dimensional subspace plus the extreme value
        theorem.  The argument that the search can be confined to a bounded
        set (by choosing any reference point $\bm{s}_0$) is correct. \checkmark
  \item \textbf{Step 2 (Uniqueness):}  Uses the parallelogram law correctly.
        The key inference is: $\|2\by - \yhat - \yhat'\|^2 \ge 4d^2$ because
        $\tfrac{1}{2}(\yhat+\yhat') \in S$, so the norm of the average
        residual is at least $d$.  Squaring and multiplying by 4 gives the
        bound. Then $\|\yhat'-\yhat\|^2 \le 0$ forces equality. \checkmark
  \item \textbf{Step 3 (Orthogonality):}  Standard variational argument.
        $f(t)$ is a quadratic in $t$ with minimum at $t=0$; setting $f'(0)=0$
        gives $\langle \by-\yhat, \bm{v}\rangle = 0$. \checkmark
  \item \textbf{Converse:}  Uses Pythagoras (cross term vanishes by
        orthogonality) to show $\|\by-\bm{s}\|^2 \ge \|\by-\yhat\|^2$.
        Correct and properly closes the biconditional. \checkmark
\end{itemize}

\textbf{No gaps, no circular arguments.}


% ----- FWL Proof -----
\subsection*{Proof of FWL (Block Elimination)}

\textbf{Verdict: CORRECT and COMPLETE.}

This is a significant improvement over the original proof sketch in
\texttt{lecture\_notes.tex} (which simply said ``solve the first block
and substitute'').  The new proof:

\begin{itemize}[nosep]
  \item Writes the block normal equations explicitly. \checkmark
  \item Solves the first block for $\bhat_1$ (valid since $\bX_1$ has
        full column rank). \checkmark
  \item Substitutes into the second block and simplifies, using the
        identity $\bX_1(\bX_1^\top\bX_1)^{-1}\bX_1^\top = \bI - \bM_1$.
        The underbrace labels are correct. \checkmark
  \item Collects terms and uses the symmetric-idempotent properties of
        $\bM_1$ to factor $\bX_2^\top\bM_1\bX_2 = (\bM_1\bX_2)^\top(\bM_1\bX_2)$
        and $\bX_2^\top\bM_1\by = (\bM_1\bX_2)^\top(\bM_1\by)$. \checkmark
  \item States the rank condition (``provided $\bM_1\bX_2$ has full column
        rank'') explicitly. \checkmark
\end{itemize}

\textbf{Minor note:}
The step ``$\bX_2^\top\bX_2 - \bX_2^\top(\bI-\bM_1)\bX_2 = \bX_2^\top\bM_1\bX_2$''
is stated in a parenthetical remark (line~162).  This is algebraically
immediate ($\bX_2^\top\bX_2 - \bX_2^\top\bX_2 + \bX_2^\top\bM_1\bX_2 =
\bX_2^\top\bM_1\bX_2$), so the omission of a detailed expansion is acceptable.

\textbf{No gaps, no circular arguments.}


% ----- Gauss-Markov Proof -----
\subsection*{Proof of Gauss--Markov}

\textbf{Verdict: CORRECT and COMPLETE.}

\begin{itemize}[nosep]
  \item \textbf{Step 1:}  Decomposes $\bm{C} = (\bX^\top\bX)^{-1}\bX^\top + \bD$.  Standard. \checkmark
  \item \textbf{Step 2:}  Derives $\bD\bX = \bm{0}$ from unbiasedness.
        The chain $E[\btilde] = \bm{C}\bX\bm{\beta} = (\bI_p + \bD\bX)\bm{\beta}$
        is correct.  Requiring this to equal $\bm{\beta}$ for all $\bm{\beta}$
        forces $\bD\bX = \bm{0}$. \checkmark
  \item \textbf{Step 3:}  Expands $\bm{C}\bm{C}^\top$ and uses $\bD\bX = \bm{0}$
        to kill the cross terms.  The key computation:
        $(\bX^\top\bX)^{-1}\bX^\top\bD^\top = (\bD\bX(\bX^\top\bX)^{-1})^\top
        = \bm{0}^\top = \bm{0}$. \checkmark

        The remaining term $(\bX^\top\bX)^{-1}\bX^\top\bX(\bX^\top\bX)^{-1}
        = (\bX^\top\bX)^{-1}$ is correct. \checkmark
  \item \textbf{Step 4:}  Concludes $\text{Var}(\btilde) - \text{Var}(\bhat)
        = \sigma^2\bD\bD^\top \succeq \bm{0}$ with equality iff $\bD = \bm{0}$.
        Correct. \checkmark
\end{itemize}

The geometric remark (``the rows of $\bD$ live in $\Col(\bX)^\perp$'')
is a correct restatement of $\bD\bX = \bm{0}$ and adds genuine intuition.

\textbf{No gaps, no circular arguments.}

\medskip
\noindent\textbf{Potential label conflict:}
The original document defines \verb|\label{thm:gm}| for the Gauss--Markov
theorem. The new proof uses \verb|\label{thm:gm-full}|. If the intention is
to \emph{replace} the original theorem statement, the label should remain
\verb|thm:gm| to avoid breaking existing references. If both are to
coexist, the labels are fine. See Part~C for details.


% ----- Characterisation Theorem -----
\subsection*{Proof of the Characterisation Theorem (Symmetric + Idempotent $\Leftrightarrow$ Orthogonal Projection)}

\textbf{Verdict: CORRECT and COMPLETE.}

\begin{itemize}[nosep]
  \item \textbf{(a) $\Rightarrow$ (b):}  Idempotency follows from
        ``projecting something already in $W$ gives itself.''  Symmetry
        follows from $\langle P\bm{u}, \bm{v}\rangle = \langle P\bm{u},
        P\bm{v}\rangle = \langle \bm{u}, P\bm{v}\rangle$, with each step
        justified by orthogonal decomposition. \checkmark
  \item \textbf{(b) $\Rightarrow$ (a):}  Part~(i) ($P\bm{v} \in W$) is
        immediate from the definition of $\Col(P)$.  Part~(ii)
        ($\bm{v} - P\bm{v} \perp W$) uses symmetry ($P^\top = P$) and
        idempotency ($P^2 = P$) in the computation
        $\bm{v}^\top P\bm{z} - \bm{v}^\top P^2\bm{z} = 0$. \checkmark
\end{itemize}

\textbf{No gaps, no circular arguments.}


% =====================================================================
%                     PART C: CONSISTENCY CHECK
% =====================================================================

\section*{Part C: Consistency with \texttt{lecture\_notes.tex}}


\subsection*{Macro usage}

All macros used in \texttt{round2\_math.tex} and \texttt{round2\_examples.tex}
are defined in the preamble of \texttt{lecture\_notes.tex}:

\begin{center}
\renewcommand{\arraystretch}{1.1}
\begin{tabular}{ll}
\toprule
Macro & Defined? \\
\midrule
\verb|\bX, \by, \bH, \bM, \be, \yhat, \bhat, \btilde| & Yes \\
\verb|\bI, \bQ, \bD| & Yes \\
\verb|\Col, \ip, \norm, \R, \tr, \rank, \diag, \argmin| & Yes \\
\verb|\bm{...}| (from \texttt{bm} package) & Yes \\
\bottomrule
\end{tabular}
\end{center}

\noindent\textbf{No undefined macros found.}

\medskip
\noindent\textbf{Potential concern:} The Gauss--Markov proof uses
\verb|\bD| for the ``excess matrix'' $\bm{C} - (\bX^\top\bX)^{-1}\bX^\top$.
The macro \verb|\bD| is already defined in the main document as
$\bm{D}$ and is used in the Ridge regression section for the diagonal
eigenvalue matrix.  Since the two uses occur in different sections and
the meaning is clear from context, this is acceptable.  However, if
there is any possibility of the Gauss--Markov proof and the Ridge
section appearing on the same page, consider using a different letter
(e.g., $\bm{A}$) for the excess matrix.


\subsection*{Label conflicts}

\begin{center}
\renewcommand{\arraystretch}{1.1}
\begin{tabular}{lll}
\toprule
New label & Existing label & Conflict? \\
\midrule
\verb|eq:fwl-blocks| & --- & No conflict \\
\verb|eq:fwl-beta1| & --- & No conflict \\
\verb|thm:gm-full| & \verb|thm:gm| & Potential issue (see below) \\
\verb|eq:gm-constraint| & --- & No conflict \\
\verb|thm:proj-char| & --- & No conflict \\
\verb|rem:hii-intercept| & --- & No conflict \\
\verb|rem:intercept-decomp| & --- & No conflict \\
\verb|sec:adj-r2| & --- & No conflict \\
\verb|sec:df| & --- & No conflict \\
\verb|ex:first-regression| & --- & No conflict \\
\verb|ex:pythagoras| & --- & No conflict \\
\verb|ex:fwl| & --- & No conflict \\
\verb|ex:multicollinearity| & --- & No conflict \\
\verb|ex:leverage| & --- & No conflict \\
\verb|ex:ridge| & --- & No conflict \\
\bottomrule
\end{tabular}
\end{center}

\noindent\textbf{Issue:} The Gauss--Markov theorem is currently labeled
\verb|thm:gm| in the original document.  The round2 version introduces a
new label \verb|thm:gm-full|.  If the intention is to \emph{replace} the
existing theorem and proof, the new label should be \verb|thm:gm| to
preserve any references.  If both are to coexist (unlikely), the new
label is fine.  \textbf{Recommendation:} Change \verb|thm:gm-full| to
\verb|thm:gm|.


\subsection*{Cross-references}

The following cross-references in the new material point to labels that
exist in \texttt{lecture\_notes.tex}:

\begin{itemize}[nosep]
  \item \verb|Proposition~\ref{prop:hat}| $\to$ \verb|\label{prop:hat}| (line~244). \checkmark
  \item \verb|Proposition~\ref{prop:r2mono}| $\to$ \verb|\label{prop:r2mono}| (line~434). \checkmark
  \item \verb|Section~\ref{sec:ftest}| $\to$ \verb|\label{sec:ftest}| (line~882). \checkmark
  \item \verb|Example~\ref{ex:first-regression}| $\to$ self-reference within round2. \checkmark
  \item \verb|Example~\ref{ex:fwl}| $\to$ self-reference within round2. \checkmark
\end{itemize}

\noindent\textbf{All cross-references resolve correctly.}


\subsection*{Theorem/example numbering}

The original document uses \verb|\newtheorem{example}[theorem]{Example}|,
which means examples share the theorem counter.  The new examples will be
numbered according to the section in which they are inserted.  Since the
insertion directives specify locations (e.g., ``after Section~2.3''),
the numbering will be assigned automatically by \LaTeX{} and will not
conflict.  \textbf{No issue.}


\subsection*{Package dependencies}

The new material uses \verb|\boxed{}| (from \texttt{amsmath}, already loaded),
\verb|\checkmark| (from \texttt{amssymb}, already loaded),
\verb|[nosep,label=(\alph*)]| and \verb|[nosep,label=(\roman*)]| (from
\texttt{enumitem}, already loaded), and \verb|\tcolorbox| (already loaded).
\textbf{No new packages needed.}


% =====================================================================
%                     PART D: CORRECTED LATEX
% =====================================================================

\section*{Part D: Corrected \LaTeX}


% =====================================================================
% CORRECTION 1: Example 3, Step B3
% =====================================================================

\subsection*{Correction 1: Example 3 (FWL), Step B3}

Replace the current Step B3 block (from ``\emph{Step B3: Residualise
$\by$ on $\bX_1$.}'' through the displayed residual $\tilde{\by}$)
with the following:

\begin{verbatim}
\emph{Step B3: Residualise $\by$ on $\bX_1$.}
\[
\bX_1^\top\by = \begin{pmatrix}12\\8\end{pmatrix},\quad
(\bX_1^\top\bX_1)^{-1}\bX_1^\top\by
= \begin{pmatrix}1/2&-1/2\\-1/2&1\end{pmatrix}
  \begin{pmatrix}12\\8\end{pmatrix}
= \begin{pmatrix}2\\2\end{pmatrix}.
\]
Projection: $\bX_1\bigl(\begin{smallmatrix}2\\2\end{smallmatrix}\bigr)
= (2+2,\;2+0,\;2+2,\;2+0)^\top = (4,\;2,\;4,\;2)^\top$.

Residual:
\[
\tilde{\by} = (2,3,6,1)^\top - (4,2,4,2)^\top
= (-2,\;1,\;2,\;-1)^\top.
\]
\end{verbatim}

And update Step B4's numerator calculation:

\begin{verbatim}
Numerator:
\[
\tilde{\bm{x}}_2^\top\tilde{\by}
= \bigl(-\tfrac{1}{2}\bigr)(-2) + \bigl(\tfrac{1}{2}\bigr)(1)
  + \bigl(\tfrac{1}{2}\bigr)(2) + \bigl(-\tfrac{1}{2}\bigr)(-1)
= 1 + \tfrac{1}{2} + 1 + \tfrac{1}{2} = 3.
\]
\end{verbatim}


% =====================================================================
% CORRECTION 2: Example 4, Case B perturbed solution
% =====================================================================

\subsection*{Correction 2: Example 4 (Multicollinearity), Case B perturbed values}

Replace the paragraph beginning ``But change $y_4$ to $10$'' with:

\begin{verbatim}
But change $y_4$ to $10$ (a shift of $+1$) and the solution becomes
approximately $\bhat \approx (3,\; -98,\; 100)$ --- the individual
coefficients swing by $\pm 100$ even though the \emph{prediction}
$\yhat$ barely changes.
\end{verbatim}

And in the summary table, replace the row:

\begin{verbatim}
$\Delta\hat\beta$ per unit $\Delta y$ & $\le 1/2$ & $\approx 67$ \\
\end{verbatim}

with:

\begin{verbatim}
$\Delta\hat\beta$ per unit $\Delta y$ & $\le 1/2$ & $\approx 100$ \\
\end{verbatim}


% =====================================================================
% CORRECTION 3 (OPTIONAL): Gauss--Markov label
% =====================================================================

\subsection*{Correction 3 (Recommended): Gauss--Markov label}

If the new Gauss--Markov theorem and proof are intended to replace the
existing ones, change:

\begin{verbatim}
\begin{theorem}[Gauss--Markov]\label{thm:gm-full}
\end{verbatim}

to:

\begin{verbatim}
\begin{theorem}[Gauss--Markov]\label{thm:gm}
\end{verbatim}


% =====================================================================
%                     SUMMARY OF FINDINGS
% =====================================================================

\section*{Summary}

\begin{center}
\renewcommand{\arraystretch}{1.3}
\begin{tabular}{llcc}
\toprule
Item & Verdict & Errors & Severity \\
\midrule
Example 1 & All correct & 0 & --- \\
Example 2 & All correct & 0 & --- \\
Example 3 & Final answer correct, intermediates wrong & 1 & ERROR \\
Example 4 & Approximate values wrong by factor $\sim 1.5$ & 1 & ERROR \\
Example 5 & All correct & 0 & --- \\
Example 6 & All correct & 0 & --- \\
\midrule
Projection Theorem proof & Complete, no gaps & 0 & --- \\
FWL proof & Complete, no gaps & 0 & --- \\
Gauss--Markov proof & Complete, no gaps & 0 & --- \\
Characterisation proof & Complete, no gaps & 0 & --- \\
\midrule
Macro consistency & No undefined macros & 0 & --- \\
Label conflicts & \verb|thm:gm-full| vs \verb|thm:gm| & 1 & NOTE \\
Cross-references & All resolve & 0 & --- \\
\bottomrule
\end{tabular}
\end{center}

\medskip
\noindent\textbf{Bottom line:} Two numerical errors need correction
(Example~3 Step~B3 and Example~4 Case~B approximate values).  All four
proofs are logically complete and free of gaps or circular reasoning.
The new material is fully consistent with the macro and labeling
conventions of the original document, subject to the minor label
recommendation above.

\noindent Corrected \LaTeX{} for both errors is provided in Part~D.
